\chapter{דו-כיווניות}

\section{ארכיטקטורת \textenglish{Sequence-to-Sequence LSTM}}

ארכיטקטורת \textenglish{Sequence-to-Sequence} המבוססת על רשתות \textenglish{LSTM} מהווה אבן דרך בתחום הלמידה העמוקה לעיבוד סדרות זמן \cite{valentini2019speech}. המודל מורכב משני רכיבים עיקריים: מקודד (\textenglish{encoder}) שמעבד את הרצף הקלט ויוצר ייצוג וקטורי מופשט, ופוענח (\textenglish{decoder}) שמשתמש בייצוג זה כדי ליצור את הרצף הפלט. בהקשר של חילוץ פרמטרים מאותות סינוסואידליים, גישה זו מאפשרת לרשת ללמוד את הדינמיקה הסיבוכה של האותות רועשים ולייצר אותם בצורה דחוסה וניתנת לעיבוד.

המקודד מעבד את ערכי הקלט ברצף ויוצר סדרה של מצבים חבויים המכילים מידע מצטבר על האות כולו. בשלב הראשון, המודל מבצע הטמעת קלט, כלומר המרת כל ערך למימד גבוה-ממדי המתאים ללמידה יעילה. לאחר מכן, המקודד \textenglish{LSTM} משדרג את המצב החבוי בהתאם לקלט הנוכחי, תוך שמירה על מידע זיכרון מטווח ארוך. בתום עיבוד כל הרצף, המקודד מוציא ``וקטור הקשר'' (\textenglish{context vector}) המהווה סיכום דחוס של המידע במקור.

הפוענח משתמש בוקטור הקשר הזה כמצב ראשוני ויוצר את הרצף הפלט צעד אחד בכל פעם. בכל שלב זמן, הפוענח מעבד את הפלט הקודם (או סימול התחלה בשלב הראשון) ויוצר את ערך הפלט הבא. מבנה זה מאפשר לרשת ללמוד קשרים מורכבים בין הקלט לפלט, בעיקר כאשר משולבים בה מנגנוני קשב מתקדמים \cite{park2024hybrid}. בטבלה הבאה מוצגים ההבדלים העיקריים בין \textenglish{RNN/LSTM} לבין ארכיטקטורת \textenglish{Transformer}:

\begin{table}[H]
\centering
\caption{השוואה בין \textenglish{RNN/LSTM} לבין \textenglish{Transformer}}
\label{tab:rnn_vs_transformer}
\begin{english}
\begin{tabular}{lll}
\toprule
\textbf{Property} & \textbf{RNN / LSTM} & \textbf{Transformer} \\
\midrule
Processing order   & Sequential (step-by-step)   & Fully parallel \\
Time complexity    & $O(n)$ serial steps         & $O(n^2 \cdot d)$ parallel \\
Long-range deps    & Difficult (vanishing grad)  & Direct (self-attention) \\
Memory footprint   & $O(1)$ hidden state         & $O(n^2)$ attention matrix \\
Training signal    & BPTT (unstable)             & Stable, parallelizable \\
Real-time inference & Natural (streaming)        & Requires full sequence \\
\bottomrule
\end{tabular}
\end{english}
\end{table}

\section{תחזוקת רגרסיה של פרמטרים: משרעת, תדר ופאזה}

הגישה של רגרסיה ישירה לתחזוקת פרמטרים של גלים סינוסואידליים מנצלת את יכולת הרשתות העמוקות ללמוד מיפויים מורכבים ממרחב הקלט למרחב הפרמטרים \cite{engel2017wavenet}. במקום לצפות מהמודל לשחזר את האות כולו, אנו מאמנים אותו לחזות בדיוק שלושה ערכים: המשרעת (\textenglish{amplitude}) \(A\), התדר (\textenglish{frequency}) \(f\), והפאזה (\textenglish{phase}) \(\phi\) של הגל הבסיסי. גישה זו מפחיתה משמעותית את מורכבות הפלט ומשפרת את הדיוק של תחזוקת הפרמטרים בתנאים רועשים.

הגדרת בעיית הרגרסיה הולכת כדלקמן: בהינתן אות \(x(t) = A\sin(2\pi f \cdot t + \phi) + n(t)\) כאשר \(n(t)\) היא רעש גאוסי, הרשת צריכה לפלוט משלוש הערכים \(\hat{A}, \hat{f}, \hat{\phi}\) כך שהשגיאה הריבועית הממוצעת (\textenglish{MSE}) בין הפרמטרים האמיתיים לחזויים מוזערת. מסקנה חשובה היא שחילוץ פרמטרים מניב מידע מובנה יותר מאשר שחזור הנקודות הגולמיות, כאשר החלק הקידודי של הרשת יכול להתמקד בתכונות ספקטרליות רלוונטיות לקביעת כל פרמטר.

השימוש בטכניקות דוגמית מונחה בעומק (\textenglish{supervised learning}) עם \textenglish{data augmentation} משיג שיפורים משמעותיים בהכללה למודלים שלא נתקלו בהם בעת האימון \cite{luo2019convtasnet}. כאשר משתמשים בנתונים סינתטיים המייצרים משרעות, תדרים ופאזות שונות בטווחים רחבים, הרשת לומדת ייצוגים חזקים המסוגלים להתמודד עם שונות רחבה. בנוסף, הפעלת טכניקות כמו \textenglish{batch normalization} וירידה מוקדמת (\textenglish{early stopping}) משפרת את היציבות של האימון ומונעת התאמת יתר.

\section{מנגנוני קשב וחיבורים שיורים}

שילוב מנגנוני קשב (\textenglish{attention mechanisms}) בתוך ארכיטקטורת \textenglish{Sequence-to-Sequence} משנה באופן יסודי את יכולת הרשת להתמקד בחלקים רלוונטיים של רצף הקלט \cite{valentini2019speech}. במקום שהפוענח יסתמך על וקטור קשר יחיד המכיל מידע מדחוס על כל הרצף, מנגנון הקשב מאפשר לו לגשת למצבים החבויים של המקודד בצורה דיפרנציאלית, עם משקלים שמתבררים במהלך האימון. לכל שלב זמן בפוענח, המודל מחשב ציוני קשב המתארים עד כמה ``קשור'' כל מצב נסתר של המקודד לעבודה הנוכחית.

הגדרת מנגנון הקשב מתבטאת באמצעות המשוואה הבאה, בה \(h_i\) הוא המצב החבוי של המקודד בשלב \(i\), ו-\(s_t\) הוא המצב החבוי של הפוענח בשלב \(t\):

\begin{equation}
\text{attention}(s_t, \{h_i\}) = \sum_{i=1}^{T} \alpha_{t,i} h_i, \quad \alpha_{t,i} = \frac{\exp(e_{t,i})}{\sum_{j=1}^{T} \exp(e_{t,j})}
\end{equation}

כאשר \(e_{t,i}\) הוא ניקוד קשב המחושב באמצעות פונקציית ניקוד (\textenglish{scoring function}) כמו מכפלה דוטית או רשת עצבית קטנה.

חיבורים שיורים (\textenglish{residual connections}) משמשים להקלה על בעיות כמו היעלמות גרדיאנטים בעומקים רבים של רשת \cite{nachmani2020voice}. בחיבור שיורי, הפלט של שכבה מחוברת ישירות לקלט שלה, כלומר \(y = f(x) + x\) במקום רק \(y = f(x)\). בהקשר של מודלים עם קשב, חיבורים שיורים מאפשרים לגרדיאנטים לזרום בקלות דרך הרשת, וגם לשמור על מידע ``בדרך כול'' המחלף בשכבות הרבות. השילוב של קשב וחיבורים שיורים יוצר ארכיטקטורה היברידית מחזקת (\textenglish{hybrid architecture}) המשיגה ביצועים עדיף על קו בסיס קלאסי בחילוץ פרמטרים של גלים סינוסואידליים \cite{tzinis2020universal}.
